% Generated draft: requires mathematical and source review.
\documentclass[11pt]{article}
\usepackage[margin=25mm]{geometry}
\usepackage{fontspec}
\setmainfont{Latin Modern Roman}
\usepackage{amsmath,amssymb,mathtools,unicode-math}
\setmathfont{Latin Modern Math}
\usepackage{enumitem,needspace,xurl,hyperref}
\setlength{\parindent}{0pt}
\setlength{\parskip}{5pt}
\setlength{\emergencystretch}{3em}
\newcommand{\N}{\mathbb{N}}
\newcommand{\Z}{\mathbb{Z}}
\newcommand{\Q}{\mathbb{Q}}
\newcommand{\R}{\mathbb{R}}
\newcommand{\C}{\mathbb{C}}
\newcommand{\F}{\mathbb{F}}
\newcommand{\A}{\mathbb{A}}
\newcommand{\PP}{\mathbb{P}}
\newcommand{\E}{\mathbb{E}}
\newcommand{\eps}{\varepsilon}
\newcommand{\dd}{\,\mathrm{d}}
\begin{document}
\section*{UnsolvedMath ID: 30003212}
\textbf{Source problem code:} catalogue entry
\subsection{Definitions and mathematical statement}
Interpret the countable group \(G\) as a discrete group, using the standard convention for the boundary theory in the supplied question; the source does not separately specify a group topology. A \(G\)-boundary is a nonempty compact Hausdorff space \(X\) with a continuous action of \(G\) that is minimal and strongly proximal. Minimality means that \(\overline{Gx}=X\) for every \(x\in X\). Strong proximality means that, for every regular Borel probability measure \(\mu\) on \(X\), the weak-* closure of \(\{g_*\mu:g\in G\}\) contains a point mass \(\delta_x\) for some \(x\in X\), where \((g_*\mu)(A)=\mu(g^{-1}A)\) for Borel sets \(A\subseteq X\).

For \(g\in G\), put \(\operatorname{Fix}_X(g)=\{x\in X:gx=x\}\). The action is topologically free if \(\operatorname{int}_X\operatorname{Fix}_X(g)=\varnothing\) for every \(g\ne e\). Thus the hypothesis in the question is
\[
\exists g\in G\setminus\{e\}\qquad \operatorname{int}_X\operatorname{Fix}_X(g)\ne\varnothing.
\]

For \(g\in G\), let \(\lambda_G(g)\) be the unitary operator on \(\ell^2(G)\) given by \((\lambda_G(g)\xi)(h)=\xi(g^{-1}h)\). The reduced group \(C^*\)-algebra is
\[
C_r^*(G)=\overline{\operatorname{span}\{\lambda_G(g):g\in G\}}^{\|\cdot\|}\subseteq\mathcal B(\ell^2(G)).
\]
The group \(G\) is \(C^*\)-simple if \(C_r^*(G)\) has no nonzero proper norm-closed two-sided ideal.

\textbf{Problem (supplied formulation).} Given a countable group \(G\) and a \(G\)-boundary \(X\) whose action is not topologically free, determine under what additional conditions \(G\) is \(C^*\)-simple.

\textbf{Unresolved scope.} The supplied formulation requests criteria but specifies neither a candidate criterion nor whether the desired conditions should concern the group, the given action, or both. It imposes no faithfulness or nontriviality requirement on the boundary action and does not identify \(X\) with the universal Furstenberg boundary. These additional restrictions are therefore not assumed. The hypothesis is failure of topological freeness, which must not be replaced by mere failure of freeness.
\subsection{Short English statement}
For a countable discrete group acting minimally and strongly proximally on a nonempty compact Hausdorff space, assume that some nonidentity element fixes a nonempty open set. The question asks for conditions ensuring simplicity of the reduced group \(C^*\)-algebra. The source leaves the form and scope of the requested criteria unspecified.
\subsection{Sources}
\begin{itemize}
\item Supplied local registry record: permanent UnsolvedMath ID \(30003212\), source problem code \texttt{catalogue entry}, in the Oberwolfach Reports Open Problems collection. The record gives proposed year 2016 and names no individual proposer.
\item Supplied report citation: \emph{Measured Group Theory (2017)}, Oberwolfach Reports, pp.\ 2347--2397, DOI \texttt{10.4171/owr/2016/41}. Source URL and literature reference labelled ``original discussion'': \url{https://doi.org/10.4171/owr/2016/41}. The report title/year and DOI year are retained as supplied; the precise problem location within the page range is not provided.
\item Supplied literature reference labelled ``Kalantar--Kennedy'': \url{https://doi.org/10.1515/crelle-2014-0070}.
\item Supplied literature reference labelled ``Le Boudec--Matte Bon'': \url{https://arxiv.org/abs/1605.01651}.
\item Literature assessment date: \texttt{2026-08-21}. This is supplied dated curation, not an independently verified assessment. The record labels the problem ``partially solved'' and states that \(C^*\)-simplicity has a complete boundary characterization and many criteria, including results concerning groups with trivial amenable radical and suitable boundary actions, while the extracted general question has no uniform classification for every group. The supplied curation expressly requires independent verification of time-sensitive status and cited claims.
\end{itemize}
\end{document}
